\subsection{QC-006: a nonzero determinant evaluation excludes the degree-six candidate}

\paragraph{Statement QC-006.}
Let $P$ be the integral highest-weight polynomial of QC-005, of
degree $6$ and weight $\lambda=(15,13,13,13)$. There is an explicit
integer $9\times9$ linear determinant pencil $L$ such that
\[
 P(\det L)=33\pmod {101}.
 \tag{QC-006.1}
\]
Consequently $P$ does not belong to the ideal of $D_{9,4}$.
Together with the ambient multiplicity $a_\lambda=1$ checked in
the main packet and the nonzero padded restriction in QC-005,
this proves $K_9(\lambda)=q_\lambda=1$. This candidate is not
a multiplicity obstruction for $(m,n)=(8,9)$.

\begin{proof}
Set
\[
 H=x_3^3+x_0x_1x_3+x_0^2x_2
 =\det\begin{pmatrix}
 x_3&x_0&0\\
 -x_1&x_3&x_0\\
 x_2&0&x_3
 \end{pmatrix},
 \qquad
 F_D=(x_0x_1x_2)^2H.
 \tag{QC-006.2}
\]
Thus $F_D=\det L$, where the explicitly specified block diagonal
pencil is
\[
 L=\operatorname{diag}\left(
 x_0,x_0,x_1,x_1,x_2,x_2,
 \begin{pmatrix}x_3&x_0&0\\-x_1&x_3&x_0\\x_2&0&x_3\end{pmatrix}
 \right).
 \tag{QC-006.3}
\]
Its three monomials have coefficient one and exponent vectors
\[
 \alpha_A=(2,2,2,3),\qquad
 \alpha_B=(3,3,2,1),\qquad
 \alpha_C=(4,2,3,0).
 \tag{QC-006.4}
\]

For completeness, the polynomial being evaluated is defined
without any choice of a highest-weight basis. Introduce six
vectors $y_0,\ldots,y_5$ with four coordinates. Write
$\Delta_{ijkl}$ for the determinant of the $4\times4$ matrix
whose columns are the coordinate derivative vectors in copies
$i,j,k,l$, in that order. Then
\[
 P(F)=\left.
 \Delta_{0135}^{4}\Delta_{1235}\Delta_{2345}^{3}
 \Delta_{0234}\Delta_{0124}^{4}
 \prod_{i=0}^{5}F(y_i)\right|_{y_4=y_5=e_0}.
 \tag{QC-006.5}
\]
There are no remaining variables in copies $0,1,2,3$.
The differentiation orders in the six copies are
$(9,9,9,9,8,8)$. The total number of derivatives in each
coordinate is $13$. Thus evaluation at $e_0$ in the last two
copies requires total input exponent $(15,13,13,13)$.
If $a,b,c$ of the six input monomials are respectively of types
$A,B,C$, the $x_2$ and $x_1$ exponents give
$12+c=13$ and $12+b=13$. Hence $a=4,b=c=1$.
There are exactly $6\cdot5=30$ ordered choices of the letter
$b$ carrying type $B$ and the distinct letter $c$ carrying type
$C$; the other four letters carry type $A$.

Here is a finite coefficient contraction that evaluates all
thirty contributions without expanding $24^{13}$ terms.
For a $4\times4$ exponent matrix $T=(T_{rs})$, define
\[
 C_k(T)=[z^T]\det(z_{rs})^k,
 \qquad 0\leq r,s\leq3,
 \tag{QC-006.6}
\]
and define it to be zero when an entry of $T$ is negative.
The table $C_k$ is obtained directly from the Leibniz expansion:
start with $C_0(0)=1$ and, $k$ times, add each permutation
matrix with its permutation sign. This is an exact recurrence
for every coefficient. Reduction modulo $101$ leaves $1848$
nonzero table entries for $k=3$ and $10147$ for $k=4$.

For a fixed pair $(b,c)$, let $\alpha_i$ be the exponent vector
assigned to letter $i$, and put
\[
 \gamma_i=\alpha_i\quad(i=0,1,2,3),\qquad
 \gamma_i=\alpha_i-e_0\quad(i=4,5).
 \tag{QC-006.7}
\]
The desired bracket coefficient is
\[
 h_{bc}=[v_0^{\gamma_0}\cdots v_5^{\gamma_5}]
 [0135]^4[1235][2345]^3[0234][0124]^4.
 \tag{QC-006.8}
\]
Rows of a bracket coefficient table are indexed by its letters,
and columns by coordinates. Transposing the usual bracket
matrix does not change its determinant.

To specify the convolution completely, choose rows
$(R_0,R_1,R_3,R_5)$ of a term of $[0135]^4$, rows
$(S_0,S_1,S_2,S_4)$ of a term of $[0124]^4$, and permutations
$\sigma,\tau\in\mathfrak S_4$ for the single brackets
$[0234]$ and $[1235]$, respectively. Retain only choices with
\[
 R_0+S_0+e_{\sigma(0)}=\gamma_0,
 \qquad
 R_1+S_1+e_{\tau(0)}=\gamma_1.
 \tag{QC-006.9}
\]
The four remaining rows required of $[2345]^3$ are
\[
 \begin{split}
 T_2&=\gamma_2-S_2-e_{\sigma(1)}-e_{\tau(1)},\\
 T_3&=\gamma_3-R_3-e_{\sigma(2)}-e_{\tau(2)},\\
 T_4&=\gamma_4-S_4-e_{\sigma(3)},\\
 T_5&=\gamma_5-R_5-e_{\tau(3)}.
 \end{split}
 \tag{QC-006.10}
\]
Then $h_{bc}$ is the sum, over these choices, of
\[
 C_4(R_0,R_1,R_3,R_5)\,
 C_4(S_0,S_1,S_2,S_4)\,
 \operatorname{sgn}(\sigma)\operatorname{sgn}(\tau)\,
 C_3(T_2,T_3,T_4,T_5).
 \tag{QC-006.11}
\]
This follows by multiplying the five bracket factors and
equating the exponent of each letter and coordinate. Every
term occurs exactly once. Grouping the terms of the second
$C_4$ table by their first two rows implements
(QC-006.9); caching the four residual rows in (QC-006.10)
only avoids repeating identical finite sums.

The standalone standard-library checker
\texttt{check\_nonic\_determinant.py} implements precisely
(QC-006.6)--(QC-006.11), with no imported symbolic algebra
system or source certificate. Its complete residue table,
with rows $b$ and columns $c$ indexed from $0$ to $5$, is
\[
 (h_{bc})=\begin{pmatrix}
 -&43&85&16&98&8\\
 43&-&16&85&8&98\\
 77&32&-&21&64&18\\
 32&77&21&-&18&64\\
 26&40&76&34&-&0\\
 40&26&34&76&0&-
 \end{pmatrix}\pmod {101}.
 \tag{QC-006.12}
\]
Its sum is $64$ modulo $101$.

For an ordinary monomial $y^\alpha$, differentiation by
$\gamma$ gives $\alpha!/(\alpha-\gamma)!$ times the
remaining monomial. Here the remainder is zero or $e_0$,
whose factorial is one. The multiplier is therefore
$\alpha_i!$ in every copy, including copies $4$ and $5$.
It is independent of $(b,c)$ and equals
\[
 (2!2!2!3!)^4(3!3!2!1!)(4!2!3!0!)
 =48^4\cdot72\cdot288
 =81\pmod {101}.
 \tag{QC-006.13}
\]
Consequently
\[
 P(F_D)=64\cdot81=33\pmod {101},
 \tag{QC-006.14}
\]
proving (QC-006.1). All data and operators have integer
coefficients, so a nonzero residue proves nonvanishing over
$\mathbb Q$ and $\mathbb C$, rather than only in the finite
field. The restriction of the indicated ambient highest-weight
line to the determinantal locus is nonzero. Ambient
multiplicity one then gives $K_9(\lambda)=1$, and QC-005
gives $q_\lambda=1$.
\end{proof}

\paragraph{Scope and verification.}
The counterexample is an actual determinant, so no closure,
generic-rank, or stabilizer upper-bound inference is involved.
The checker reconstructs its own coefficient tables and all
thirty contributions, and asserts the term counts, sum,
derivative factor, and final value. The stored output records
every summand as well as the final residue. The proof and
checker settle this specific candidate; they do not establish
containment of the whole padded variety in $D_{9,4}$ or exclude
other weights.

\paragraph{Geometric restatement.}
The equivariant Schur-functor section constructed in QC-005
has nonzero restriction at the integer point (QC-006.3) of the
determinantal parameter image. Thus its line is not a submodule
of the determinantal ideal sheaf. The failure of this proposed
obstruction is witnessed by a section value, independently of
any dimension count. The grading and Schur-operation framework
is the same SGA~3/SGA~6 dictionary as QC-005.
